\documentclass[10pt,a4paper]{amsart}
\usepackage[margin=19mm]{geometry}
\usepackage{amsmath,amssymb,mathtools}
\usepackage[scr=rsfs,cal=euler]{mathalfa}
\usepackage{xfrac}
\usepackage[unicode,hidelinks,bookmarks=false]{hyperref}
\newcommand{\ie}{i.e.\ }
\newcommand{\cf}{cf.\ }
\newcommand{\resp}{respectively\ }
\newcommand{\Subsection}{\subsection}
\providecommand{\nobreakdash}{\nobreak\hbox{-}}
\setlength{\emergencystretch}{2em}
\multlinegap0pt
\usepackage{amssymb}
\usepackage{mathtools}
\usepackage[shortlabels]{enumitem}
\usepackage{xcolor}
\usepackage{float}
\usepackage{tikz}
\newtheorem{lemma}{Lemma}
\title{Zeta Renormalization and Pressure at Infinity\\ for an Infinitely Cusped Tree Lattice}
\author{Sanghoon Kwon}
\date{Source: arXiv:2608.25786v1 (CC BY 4.0)}
\begin{document}
\maketitle

\noindent\emph{Source and licence.} Zeta Renormalization and Pressure at Infinity\\ for an Infinitely Cusped Tree Lattice. Version arXiv:2608.25786v1 (CC BY 4.0), distributed by arXiv under the Creative Commons Attribution 4.0 International licence, http://creativecommons.org/licenses/by/4.0/, as recorded in the arXiv licence metadata for this exact version and shown on the abstract page https://arxiv.org/abs/2608.25786v1. Full licence notice: \texttt{licenses/arxiv-2608.25786-CC-BY-4.0.txt}.\par\smallskip

\noindent\textit{Editorial scope.} Counted unit: the article's lemma lem:finite-core-escape (source0.tex lines 1516--1528). The article's complete original proof of it is delivered with it (source0.tex lines 1530--1585, one \texttt{proof} environment of 56 lines). No mathematics is written, repaired or re-derived: the statement and the proof are the article's own lines, in the article's own notation, and no retained line is substituted or re-flowed.  Retained uncounted as the source-local context the proof needs: lines 875--884; lines 1082--1084.  Nothing was omitted from any retained span, so the package carries no omission record.  No retained line contains a citation, so the wrapper carries no bibliography and no bibliography entry.  No retained line contains a figure environment, an image-inclusion command or drawing code, so no image is delivered and the package declares no asset dependency.  Editorial formatting only, in addition: an AMS \texttt{amsart} preamble that restates the article's own declarations and macros used by the retained lines, namely the environments \texttt{lemma} on separate counters, together with the standard packages those lines need.  The retained environments print in this document as Lemma 1 in order of appearance, whereas the article prints the same items with the numbering of its own full document.  The complete native source of the article as distributed in the arXiv e-print for this exact version is bundled unchanged as \texttt{sources/208597/source0.tex}.
\begin{align}
 \operatorname{Per}_a(n;F,M)
 &=\left\{x\in\operatorname{Per}_a(n):
 \#\{0\le j<n:x_j\in F\}\le\frac nM\right\},
 \label{eq:escaping-periodic-set}\\
 P_\infty(\phi;a,F,M)
 &=\limsup_{n\to\infty}\frac1n
 \log\sum_{x\in\operatorname{Per}_a(n;F,M)}e^{S_n\phi(x)},
 \label{eq:escaping-periodic-pressure}
\end{align}

\begin{equation}\label{eq:excursionweight}
 W_{\rm vert}(k)=q(q+1)^{k-1}.
\end{equation}

\begin{lemma}[Finite-core escape]
\label{lem:finite-core-escape}
For the locally finite squared root component and the bounded locally constant
potential $\phi_q$,
\begin{equation}\label{eq:finite-core-pressure-limit}
 P_\infty(\phi_q)
 =\lim_{N\to\infty}
 P_G\!\left(\phi_q\bigm|
 \Sigma^{(2),+}_{a_0}\setminus K_N\right).
\end{equation}
On the right, the pressure of a nontransitive restriction is the supremum of
the pressures of its recurrent components.
\end{lemma}

\begin{proof}
We give the excursion estimate specialized to the present locally finite
graph.  Let $P_N$ denote the pressure of the unique recurrent component of
the complement of $K_N$.  For every $\delta>0$, the total weight of length
$m$ paths between any two states in the finite boundary of that component is
at most
\begin{equation}\label{eq:open-tail-bound}
 C_\delta e^{(P_N+\delta)m}.
\end{equation}
To see this uniformly, for each ordered pair $(s,t)$ of boundary states choose
one positive path $\gamma_{t,s}$ from $t$ back to $s$.  Concatenation
$\pi\mapsto\pi\gamma_{t,s}$ injects the length-$m$ paths from $s$ to $t$ into
the based periodic paths at $s$.  The added lengths and the reciprocals of the
added weights range over finite sets because the boundary is finite.  The
definition of $P_N$ therefore bounds all sufficiently long nonzero periodic
partition sums by $e^{(P_N+\delta)m}$ up to one uniform multiplicative
constant; enlarging that constant handles the finitely many short lengths.
A severed vertical-ray excursion has two-step growth at most $q+1<4q$, by the
explicit weights in \eqref{eq:excursionweight}, so after one further increase
of $C_\delta$ the same estimate applies to every component of the complement.

A periodic path counted in \eqref{eq:escaping-periodic-pressure} decomposes at
its visits to $K_N$ into at most $s\le n/M+1$ outside excursions and joining
pieces in $K_N$.  If $k\le n/M$ is the total time spent in $K_N$, finiteness
of $K_N$ and its boundary bounds the total weight of all joining pieces with
specified endpoints by $D^{k+s}$ for a constant $D$.  The factors
$C_\delta^sD^{k+s}$, together with the choices of excursion lengths and visit
positions, contribute, for a constant $C_1>0$ independent of $n$ and $M$, at most
\[
 D^{n/M}\sum_{s\le n/M+1}\binom ns C_1^s
 =\exp\!\bigl(n\,o_M(1)\bigr),
 \qquad o_M(1)=O\!\left(\frac{1+\log M}{M}\right)
 \longrightarrow0\quad(M\to\infty).
\]
Combining this with \eqref{eq:open-tail-bound} gives
\[
 P_\infty(\phi_q;a_0,K_N,M)
 \le P_N+\delta+o_M(1).
\]
First let $M\to\infty$ and then $\delta\downarrow0$.  This proves the upper
bound in \eqref{eq:finite-core-pressure-limit}.

Conversely, fix a finite state set $F$.  Choose a homogeneous tail state $e$
beyond every horizontal cell meeting $F$, in the parity component of $a_0$.
There are fixed positive paths from $a_0$ to $e$ and from $e$ back to $a_0$;
write $L$ and $c>0$ for their total length and weight.  Appending these paths
to every length-$m$ loop based at $e$ and contained in the cut homogeneous
tail component produces distinct root-based loops of length $m+L$ and total
weight at least $c$ times the cut-tail partition sum.  The middle loop avoids
$F$, so the resulting root-based loop visits $F$ at most $L$ times.  For each
fixed $M$ it therefore satisfies the escape condition once $m\ge ML$.
Taking the limsup along lengths realizing the cut-tail pressure proves
$P_\infty(\phi_q)\ge P_N$.  Translation conjugates the two parity copies of
the homogeneous tail, so $P_N$ is independent of $N$ and the asserted limit
follows.
\end{proof}

\end{document}
